% ===== 全局 Markdown -> PDF 头文件（xelatex）=====
\usepackage{graphicx}
\usepackage{unicode-math}          % 同时加载 fontspec
\defaultfontfeatures{Scale=MatchLowercase}
\defaultfontfeatures[\rmfamily]{Ligatures=TeX,Scale=1}

% 字体：等价于 -V mainfont / -V monofont / -V CJKmainfont
% 用 IfFontExistsTF 做兜底，换机器（Linux）也不会直接编译失败
\IfFontExistsTF{Inter}
  {\setmainfont{Inter}}
  {\setmainfont{Noto Sans}}
\IfFontExistsTF{JetBrainsMono Nerd Font}
  {\setmonofont{JetBrainsMono Nerd Font}}
  {\setmonofont{DejaVu Sans Mono}}
\usepackage[space]{xeCJK}          % space = 中英自动间距（pandoc 设 CJKmainfont 时默认就是它）
\IfFontExistsTF{PingFang SC}
  {\setCJKmainfont{PingFang SC}\setCJKmonofont{PingFang SC}}
  {\setCJKmainfont{Noto Sans CJK SC}}

% 页面（原来 YAML 里的 papersize + geometry）
\usepackage[a4paper,margin=2.5cm]{geometry}

% 原来 preamble.tex 的内容
\usepackage{indentfirst}
\setlength{\parindent}{2em}
\let\pandocoldincludegraphics\includegraphics
\RenewDocumentCommand{\includegraphics}{s O{} m}{%
  \noindent
  \IfBooleanTF{#1}
    {\pandocoldincludegraphics*[#2]{#3}}
    {\pandocoldincludegraphics[#2]{#3}}%
}
